\documentclass[12pt,a4paper]{article}
\usepackage[utf8]{inputenc}
\usepackage[spanish]{babel}
\usepackage{amsmath}
\usepackage{amsfonts}
\usepackage{amssymb}
\usepackage{graphicx}
\usepackage{hyperref}
\usepackage{geometry}
\geometry{a4paper, margin=2.5cm}
\usepackage{titlesec}
\usepackage{xcolor}
\usepackage{listings}

\title{\textbf{Reporte de Avances Técnicos y Científicos}\\ \large Prototipo para el análisis automatizado del comportamiento en modelos de ansiedad utilizando IA (TT 2026-A155)}
\author{Equipo de Investigación TT}
\date{\today}

\begin{document}

\maketitle

\begin{abstract}
El presente reporte documenta los avances recientes y mejoras estructurales implementadas en el sistema de fenotipado conductual en roedores (EPM). Las mejoras abarcan el rediseño de la heurística espacial para la detección de comportamientos complejos (Thigmotaxis), la optimización de la persistencia de datos y el auto-despliegue del entorno, así como el refinamiento sustancial de la interfaz de usuario en Streamlit para el entorno clínico y de laboratorio.
\end{abstract}

\section{Introducción}
El análisis continuo del laberinto en cruz elevado (EPM) presenta retos de alta precisión debido a la oclusión, la similitud visual del sujeto con el entorno y la definición espacial de variables conductuales complejas como el \textit{thigmotaxis} (tendencia a permanecer cerca de las paredes). Las recientes iteraciones tecnológicas en este proyecto han migrado heurísticas de estimación simple a análisis geométricos y algorítmicos robustos.

\section{Avances y Mejoras Recientes}

\subsection{Segmentación Avanzada de Paredes (Muros)}
Anteriormente, los bordes físicos del laberinto se inferían únicamente a partir de la superposición de \textit{bounding boxes} rectangulares, lo que presentaba limitantes para identificar interacciones reales del ratón con las barreras de acrílico. Se ha introducido un módulo de \textbf{Dibujo Vectorial Lineal} en la configuración de ROI (Region of Interest) a través de \texttt{streamlit-drawable-canvas}. El usuario experimental puede ahora trazar vectores rectos (\textit{Líneas Cyan}) para delimitar exactamente las fronteras físicas de la arena, diferenciándolas de las regiones espaciales.

\subsection{Algoritmia de Thigmotaxis Geométrica}
La detección de Thigmotaxis se ha optimizado usando la menor distancia Euclidiana de la nariz del ratón hacia los segmentos de pared definidos por el usuario. La lógica base de la función \texttt{distancia\_punto\_segmento} calcula la proyección ortogonal del punto $(x_0, y_0)$ sobre la línea delineada $(x_1, y_1)$ a $(x_2, y_2)$.
\begin{enumerate}
    \item Si el usuario trazó muros, el algoritmo valida matemáticamente la cercanía sub-píxel con margen dinámico (22 px).
    \item Se anula el registro del muro como "Zona habitable", previniendo bitácoras corruptas.
\end{enumerate}

\subsection{Ingesta y Extracción de Datos Sólida}
Se resolvieron vulnerabilidades de API causadas por anidación prohibida en Streamlit (\texttt{StreamlitAPIException}) que corrompían el registro de videos. Se implementó un \textit{Select Box dinámico} híbrido para el autocompletado de \textbf{Tratamientos} basado en los registros históricos de la base de datos PostgreSQL, mejorando la integridad referencial y facilitando los ensayos a ciegas.

\subsection{Automatización de Despliegue (Docker)}
Para mejorar la reproducibilidad en laboratorios con limitados conocimientos técnicos, se modificaron las políticas del orquestador \texttt{docker-compose.yml} agregando la regla \texttt{restart: unless-stopped}. Las bases de datos se levantan ahora silenciosamente y en \textit{background} durante el arranque del sistema operativo Windows, junto con los ejecutables \texttt{.bat} silenciados, abstraiendo totalmente el manejo del backend al investigador.

\subsection{Interfaz Analítica Refinada}
La interfaz de presentación de Historial Clínico (\texttt{Resultados y Estadísticas}) se actualizó a un \textit{layout=wide} nativo que abarca toda la pantalla. El espacio negativo ha sido reducido, integrando tarjetas de KPIs responsivas y permitiendo edición metadatos *in-line* con actualizaciones retroactivas sin romper relaciones de llaves foráneas en \texttt{SQLAlchemy}.

\section{Conclusión y Siguientes Pasos}
Las implementaciones documentadas elevan el prototipo de prueba de concepto a una herramienta estable para la fase experimental formal. La extracción de Keypoints a través de DeepLabCut / YOLO ahora goza de soporte matemático depurado. Los siguientes pasos sugieren la ejecución ininterrumpida de validaciones con videos del ensayo ciego para solidificar las métricas de varianza.

\end{document}
